% Sección: Parte 4 – Demostración práctica

% Frame 1: Presentación de ejemplos
\begin{frame}{Demostración práctica}
  \begin{itemize}
    \item Resolvemos dos problemas vistos en la Parte 1:
      \begin{enumerate}
        \item Ecuación de Poisson en 1D (u''(x) = -\sin(\pi x))
        \item Ecuación de ondas simplificada (u_{tt} = c^2 u_{xx})
      \end{enumerate}
    \item Usamos PINNs via DeepXDE para aproximar la solución.
  \end{itemize}
\end{frame}

% Frame 2: Ejemplo 1 – Poisson 1D
\begin{frame}[fragile]{Ejemplo 1: Poisson 1D}
  Código principal:
  \begin{block}{Python}
    \begin{verbatim}
import deepxde as dde
import numpy as np

def pde_poisson(x, y):
    return dde.grad.hessian(y, x) + np.sin(np.pi * x)

def bc_poisson(_, on_boundary):
    return on_boundary


    \end{verbatim}
  \end{block}
\end{frame}

\begin{frame}[fragile]{Ejemplo 1: Poisson 1D}
  Código principal:
  \begin{block}{Python}
    \begin{verbatim}
domain = dde.geometry.Interval(0, 1)
bc1 = dde.DirichletBC(domain, lambda x: 0, bc_poisson)
data1 = dde.data.PDE(domain, pde_poisson, bc1, num_domain=100, num_boundary=2)
net1 = dde.maps.FNN([1] + [50]*3 + [1], "tanh", "Glorot normal")
model1 = dde.Model(data1, net1)
model1.compile("adam", lr=0.001)
losshistory1, _ = model1.train(epochs=5000)
    \end{verbatim}
  \end{block}
\end{frame}

% Frame 3: Ejemplo 2 – Ecuación de ondas simplificada
\begin{frame}[fragile]{Ejemplo 2: Ecuación de ondas}
  Código principal:
  \begin{block}{Python}
    \begin{verbatim}
import deepxde as dde
import numpy as np

def wave_pde(x, y):  # x: [t, x]
    d2y_tt = dde.grad.hessian(y, x, (0, 0))
    d2y_xx = dde.grad.hessian(y, x, (1, 1))
    c = 1.0
    return d2y_tt - c**2 * d2y_xx

def ic(x):
    return np.sin(np.pi * x[:, 1:])

    \end{verbatim}
  \end{block}
\end{frame}

\begin{frame}[fragile]{Ejemplo 2: Ecuación de ondas}
  Código principal:
  \begin{block}{Python}
    \begin{verbatim}

def bc_wave(x, on_boundary):
    return on_boundary

geo = dde.geometry.TimeDomain(0, 1) * dde.geometry.Interval(0, 1)
ic1 = dde.IC(geo, ic, lambda x, _: np.isclose(x[0], 0))
bc2 = dde.DirichletBC(geo, lambda x: 0, bc_wave)
data2 = dde.data.TimePDE(geo, wave_pde, [ic1, bc2], num_domain=2000, num_initial=100, num_boundary=100)
net2 = dde.maps.FNN([2] + [50]*4 + [1], "tanh", "Glorot normal")
model2 = dde.Model(data2, net2)
model2.compile("adam", lr=0.001)
losshistory2, _ = model2.train(epochs=8000)
    \end{verbatim}
  \end{block}
\end{frame}

% Frame 4: Resultados y discusión
\begin{frame}{Resultados y discusión}
  \begin{itemize}
    \item Comparación gráfica de la solución aproximada vs. solución analítica (si existe).
    \item Convergencia del error durante el entrenamiento.
    \item Observaciones: precisión, costo computacional y restricciones.
  \end{itemize}
\end{frame}

